% ============================================================================
% Section D: Analyze the Concurrent Life Cycles (Deliverable D)
% ============================================================================
\section{Concurrent Life-Cycle Analysis}
\label{sec:D}

Delivering UA-EWRRF requires managing \textbf{four concurrent life cycles} in parallel from conceptual design onward (Figure~\ref{fig:concurrent_lc}): (1)~the operational recovery plant, (2)~construction and fabrication capability, (3)~support and maintenance infrastructure, and (4)~retirement and site-closure capability. Table~\ref{tab:annexD} summarizes their dependencies and late-development risks (Annexure~D).

% ============================================================================
% Figure 3: Four Concurrent Life Cycles Diagram (Deliverable D)
% ============================================================================
\begin{figure}[htbp]
\centering
\resizebox{\textwidth}{!}{%
\begin{tikzpicture}

% --- Horizontal Timeline Grid ---
\draw[->, line width=0.8pt, draw=black!80] (2.9, 8.65) -- (17.2, 8.65) node[right, font=\rmfamily\scriptsize\bfseries] {Time};
\foreach \x/\lbl in {3.4/Need, 5.4/Conceptual, 7.5/Prelim., 9.5/Detailed, 11.5/Construct, 13.8/Operation (Y1--Y15), 16.3/Retirement} {
    \draw[dashed, draw=RuleGray] (\x, 1.35) -- (\x, 8.65);
    \node[font=\rmfamily\tiny\bfseries\color{black!85}, above] at (\x, 8.68) {\lbl};
}

% --- Swimlane 1: Product / Operational System Life Cycle (y = 7.6) ---
\node[sebanner, fill=DiagBlueFill, draw=DiagBlueDraw, text width=2.45cm, minimum height=1.22cm]
    (l1) at (1.35, 7.6) {1. Operational\\System LC\\{\tiny(UA-EWRRF)}};

\node[sebox, text width=1.45cm, minimum height=0.88cm] (p_need) at (3.95, 7.6) {Need \&\\EPR Scope};
\node[sebox, text width=2.45cm, minimum height=0.88cm] (p_des)  at (6.75, 7.6) {Concept, Prelim.\\\& Detailed Design};
\node[sebox, text width=2.05cm, minimum height=0.88cm] (p_bld)  at (10.15, 7.6) {Build, Install\\\& CTO Test};
\node[seboxTeal, text width=3.15cm, minimum height=0.88cm] (p_ops) at (13.65, 7.6) {\textbf{Capability Ready (Y1)}\\10 t/d Collection, Recovery\\\& Mid-Life Upgrading};
\node[seboxAmber, text width=1.45cm, minimum height=0.88cm] (p_ret) at (16.55, 7.6) {Phase-Out\\(Year 15)};

\draw[searrow] (p_need) -- (p_des);
\draw[searrow] (p_des) -- (p_bld);
\draw[searrow] (p_bld) -- (p_ops);
\draw[searrow] (p_ops) -- (p_ret);

% --- Swimlane 2: Production / Construction Capability Life Cycle (y = 5.6) ---
\node[sebanner, fill=DiagTealFill, draw=DiagTealDraw, text width=2.45cm, minimum height=1.22cm]
    (l2) at (1.35, 5.6) {2. Construction\\Capability LC\\{\tiny(Civil \& Skids)}};

\node[seboxTeal, text width=2.85cm, minimum height=0.88cm] (c_plan) at (6.75, 5.6) {\textbf{Plan Starts: Concept}\\EPC Vendor, MIDC Plot,\\Skid Fabrication Prep};
\node[seboxTeal, text width=2.75cm, minimum height=0.88cm] (c_exec) at (10.35, 5.6) {\textbf{Capability Ready (Y$-1$)}\\Epoxy Slab, MEP, HEPA,\\Skid Commissioning};
\node[sebox, text width=2.65cm, minimum height=0.88cm] (c_mod)  at (13.95, 5.6) {Mid-Life Bay Expansion\\(Li-ion / NIR Skid Add-on)};

\draw[searrow] (c_plan) -- (c_exec);
\draw[searrow] (c_exec) -- (c_mod);

% --- Swimlane 3: Support & Maintenance System Life Cycle (y = 3.6) ---
\node[sebanner, fill=LightTint, draw=black!60, text width=2.45cm, minimum height=1.22cm]
    (l3) at (1.35, 3.6) {3. Support \&\\Maint.\ LC\\{\tiny(Spares/EHS/IT)}};

\node[sebox, text width=2.95cm, minimum height=0.88cm] (s_plan) at (8.35, 3.6) {\textbf{Plan Starts: Prelim.}\\Spares Stocking, EHS Lab,\\CMMS \& Tech Training};
\node[seboxTeal, text width=3.95cm, minimum height=0.88cm] (s_ops) at (13.15, 3.6) {\textbf{Capability Ready (Pre-Y1)}\\Blade Swap (MTTR $\le 4$ h), XRT Calib.,\\EPR API \& Worker Biomonitoring};
\node[seboxAmber, text width=1.65cm, minimum height=0.88cm] (s_end) at (16.55, 3.6) {Spares Close\\\& Support End};

\draw[searrow] (s_plan) -- (s_ops);
\draw[searrow] (s_ops) -- (s_end);

% --- Swimlane 4: Retirement, Recovery & Disposal System Life Cycle (y = 1.8) ---
\node[sebanner, fill=DiagRedFill, draw=DiagRedDraw, text width=2.45cm, minimum height=1.12cm]
    (l4) at (1.35, 1.8) {4. Retirement\\\& Disposal LC\\{\tiny(Site Closure)}};

\node[seboxAmber, text width=2.95cm, minimum height=0.88cm] (r_plan) at (6.75, 1.8) {\textbf{Plan Starts: Concept}\\Bolted Anchors, Escrow,\\TSDF \& Closure Spec};
\node[sebox, text width=3.95cm, minimum height=0.88cm] (r_esc)  at (11.55, 1.8) {Annual Closure Escrow Accrual,\\Soil Baseline Audits \& Skid Asset Log};
\node[seboxRisk, text width=2.65cm, minimum height=0.88cm] (r_exec) at (16.05, 1.8) {\textbf{Capability Ready (Y15)}\\Decontam., Skid Resale,\\MPCB Site Closure NOC};

\draw[searrow] (r_plan) -- (r_esc);
\draw[searrow] (r_esc) -- (r_exec);

% --- Clean Cross-Cycle Dependencies & Information Exchange Arrows ---
\draw[sefeedback] (p_des.south) -- node[selabel] {Plant Layout / DDPs} (c_plan.north);
\draw[sefeedback] (c_exec.north) -- node[selabel] {Commissioned Plant} (p_bld.south);
\draw[sefeedback, draw=DiagTealDraw] ($(p_des.south east)+(-0.1,0)$) -- ++(0,-0.45) -| node[selabel, pos=0.86] {MTTR / BOM} (s_plan.north);
\draw[sefeedback, draw=DiagTealDraw] ($(s_ops.north east)+(-0.15,0)$) -- ++(0,0.48) -| (15.52, 6.6) |- node[selabel, pos=0.25] {Availability $A_o$} ($(p_ops.east)+(0,-0.18)$);
\draw[sefeedback, draw=DiagOrgDraw] (r_plan.west) -| node[selabel, pos=0.62] {De-contam.\ Reqs (DDP-7)} (4.72, 6.6) |- ($(p_des.west)+(0,-0.15)$);
\draw[seriskarrow] ($(p_ret.south)+(-0.35,0)$) -- ++(0,-0.35) -| (15.58, 2.75) -- node[selabel, pos=0.45] {Y15 Handover} (15.58, 2.26);

% --- Dedicated Bottom Strip: Consequence of Late Development Callouts ---
\node[seboxRisk, text width=5.2cm, minimum height=0.85cm, font=\rmfamily\tiny, anchor=west] at (0.15, 0.35)
    {\textbf{If Late (LC-2 Construction):} Permanent grouted slabs delay CTO by 6--9 mo and physically block Year-6 modular bay retrofits.};

\node[seboxRisk, text width=5.3cm, minimum height=0.85cm, font=\rmfamily\tiny, anchor=west] at (5.85, 0.35)
    {\textbf{If Late (LC-3 Support/Maint.):} Imported blade/XRT lead times cut $A_o < 75\%$; uncalibrated sensors cut MRR $<80\%$ \& risk CTO suspension.};

\node[seboxRisk, text width=5.55cm, minimum height=0.85cm, font=\rmfamily\tiny, anchor=west] at (11.65, 0.35)
    {\textbf{If Late (LC-4 Retirement):} Unfunded heavy-metal soil remediation ($\ge$INR 2.5--4.0 Cr), zero skid resale, and MPCB site-closure refusal.};

\end{tikzpicture}%
}
\caption{Concurrent Life-Cycle Diagram for UA-EWRRF Showing Planning Initiation, Capability Readiness, Inter-Cycle Dependencies, and Late-Consideration Consequences.}
\label{fig:concurrent_lc}
\end{figure}


\begingroup
\footnotesize
\renewcommand{\arraystretch}{1.22}
\begin{xltabular}{\textwidth}{L{2.3cm} Y L{2.0cm} Y Y}
\caption{Analysis of the four concurrent life cycles of UA-EWRRF (Annexure D format).}
\label{tab:annexD}\\
\toprule
\hdrrow \tblhead{Life cycle} & \tblhead{What must be created or prepared?} & \tblhead{When should planning begin?} & \tblhead{Dependence on other life cycles} & \tblhead{Consequence of late consideration} \\
\midrule
\endfirsthead
\multicolumn{5}{l}{\footnotesize\color{SlateMuted}\textit{Table~\ref{tab:annexD} continued from previous page.}}\\
\toprule
\hdrrow \tblhead{Life cycle} & \tblhead{What must be created or prepared?} & \tblhead{When should planning begin?} & \tblhead{Dependence on other life cycles} & \tblhead{Consequence of late consideration} \\
\midrule
\endhead
\bottomrule
\endfoot

\textbf{1. Operational System LC} & 10 t/d plant, sorting skids, HEPA/ZLD, 12 kiosks, 4 vans, EPR portal. & Phase 1 (Need). & Depends on LC-2 build, LC-3 availability ($A_o \ge 88\%$), and LC-4 closure rules. & Unmet EPR quotas, safety incidents, premature obsolescence. \\
\midrule
\textbf{2. Construction Capability LC} & MIDC plot, epoxy slabs, crane rigging, skid jigs, SAT test rigs. & Phase 2 (Concept). & Uses LC-1 DDPs; embeds LC-3 hoist clearances and LC-4 bolted anchors. & 6--9 mo delay; grouted slabs block Year-6 modular upgrades. \\
\midrule
\textbf{3. Support \& Maint.\ LC} & Wear-blade stock, XRT calibration, CMMS, EHS lab, DISH training. & Phase 3 (Prelim.). & Uses LC-1 BOM/MTTR specs; feeds reliability data back to LC-1. & Spare delays cut $A_o < 75\%$; uncalibrated sensors cut MRR $<80\%$. \\
\midrule
\textbf{4. Retirement \& Disposal LC} & Baseline soil survey, closure escrow, TSDF quota, MPCB NOC plan. & Phase 2 (Concept); survey pre-Ph.\ 5. & Imposes DDP-7 constraints on LC-1/LC-2; uses LC-3 incident logs. & Unfunded soil remediation ($\ge$INR 2.5 Cr) \& MPCB closure denial. \\

\end{xltabular}
\endgroup

\subsection{Why Support, Construction, and Retirement Cannot Be Deferred}
\begin{enumerate}
    \item \textbf{Physical Lock-In:} Overhead hoist access, quick-change blade housings (DDP-3), and sealed epoxy containment with bolted skid rails (DDP-7) must be designed before pouring concrete in Phase~5; retrofitting them at Year~15 is physically impossible.
    \item \textbf{Lead-Time Synchronization:} Eddy-current rotors, XRT lamps, and DISH-certified technicians require 4--9 months of preparation before Day-1 startup.
    \item \textbf{Baseline \& Escrow Continuity:} Without a pre-construction soil baseline and Year-1 closure escrow accrual, the promoter cannot legally separate prior industrial contamination from plant operations at Year~15.
\end{enumerate}
